% math0 manuscript. Build with: latexmk -pdf paper.tex
\documentclass[11pt]{amsart}
\usepackage[margin=1.1in]{geometry}
\usepackage[T1]{fontenc}
\usepackage{lmodern}
\usepackage{amsmath,amssymb,amsthm,mathtools,mathrsfs}
\usepackage{url}
\usepackage{hyperref}
\usepackage{xcolor}
\usepackage{listings}
\usepackage{newunicodechar}
\usepackage{booktabs}

\lstdefinelanguage{lean}{
  keywords={theorem,lemma,def,abbrev,example,fun,by,have,show,from,let,in,if,then,else,match,with,
    namespace,end,open,import,section,variable,instance,structure,where,calc,obtain,intro,exact,
    refine,rcases,rintro,simp,omega,decide,norm_num,linarith,nlinarith,positivity,induction,cases,
    private,unfold,rw,at,infer_instance},
  sensitive=true, comment=[l]{--}, morecomment=[s]{/-}{-/}, string=[b]"}
\lstset{language=lean, basicstyle=\ttfamily\footnotesize, columns=fullflexible, keepspaces=true, breaklines=true,
  keywordstyle=\color{blue!60!black}\bfseries, commentstyle=\color{gray}, frame=single,
  extendedchars=true, inputencoding=utf8,
  literate={ℕ}{{$\mathbb{N}$}}1 {∀}{{$\forall$}}1 {∃}{{$\exists$}}1 {∑}{{$\sum$}}1
    {∈}{{$\in$}}1 {∧}{{$\wedge$}}1 {≤}{{$\le$}}1 {→}{{$\to$}}1 {¬}{{$\neg$}}1
    {·}{{$\cdot$}}1}
% Unicode commonly found in Lean sources.
\newunicodechar{ℕ}{\ensuremath{\mathbb{N}}}
\newunicodechar{ℤ}{\ensuremath{\mathbb{Z}}}
\newunicodechar{ℚ}{\ensuremath{\mathbb{Q}}}
\newunicodechar{ℝ}{\ensuremath{\mathbb{R}}}
\newunicodechar{ℂ}{\ensuremath{\mathbb{C}}}
\newunicodechar{∀}{\ensuremath{\forall}}
\newunicodechar{∃}{\ensuremath{\exists}}
\newunicodechar{→}{\ensuremath{\to}}
\newunicodechar{←}{\ensuremath{\leftarrow}}
\newunicodechar{↔}{\ensuremath{\leftrightarrow}}
\newunicodechar{↦}{\ensuremath{\mapsto}}
\newunicodechar{¬}{\ensuremath{\neg}}
\newunicodechar{∧}{\ensuremath{\wedge}}
\newunicodechar{∨}{\ensuremath{\vee}}
\newunicodechar{≤}{\ensuremath{\le}}
\newunicodechar{≥}{\ensuremath{\ge}}
\newunicodechar{≠}{\ensuremath{\ne}}
\newunicodechar{∈}{\ensuremath{\in}}
\newunicodechar{∉}{\ensuremath{\notin}}
\newunicodechar{⊆}{\ensuremath{\subseteq}}
\newunicodechar{∩}{\ensuremath{\cap}}
\newunicodechar{∪}{\ensuremath{\cup}}
\newunicodechar{∑}{\ensuremath{\sum}}
\newunicodechar{∏}{\ensuremath{\prod}}
\newunicodechar{⟨}{\ensuremath{\langle}}
\newunicodechar{⟩}{\ensuremath{\rangle}}
\newunicodechar{·}{\ensuremath{\cdot}}
\newunicodechar{∘}{\ensuremath{\circ}}
\newunicodechar{⁻}{\ensuremath{^{-}}}
\newunicodechar{¹}{\ensuremath{^{1}}}
\newunicodechar{²}{\ensuremath{^{2}}}
\newunicodechar{₀}{\ensuremath{_{0}}}
\newunicodechar{₁}{\ensuremath{_{1}}}
\newunicodechar{₂}{\ensuremath{_{2}}}
\newunicodechar{α}{\ensuremath{\alpha}}
\newunicodechar{β}{\ensuremath{\beta}}
\newunicodechar{ε}{\ensuremath{\varepsilon}}
\newunicodechar{δ}{\ensuremath{\delta}}
\newunicodechar{π}{\ensuremath{\pi}}
\newunicodechar{σ}{\ensuremath{\sigma}}
\newunicodechar{φ}{\ensuremath{\varphi}}
\newunicodechar{∞}{\ensuremath{\infty}}
\newunicodechar{⌊}{\ensuremath{\lfloor}}
\newunicodechar{⌋}{\ensuremath{\rfloor}}
\newunicodechar{⌈}{\ensuremath{\lceil}}
\newunicodechar{⌉}{\ensuremath{\rceil}}
\newunicodechar{∣}{\ensuremath{\mid}}
\newunicodechar{≡}{\ensuremath{\equiv}}
\newunicodechar{×}{\ensuremath{\times}}
\newunicodechar{⧸}{\ensuremath{/}}
\newunicodechar{‖}{\ensuremath{\|}}
\newunicodechar{√}{\ensuremath{\surd}}
\newunicodechar{∅}{\ensuremath{\emptyset}}
\newunicodechar{⦃}{\ensuremath{\{\!|}}
\newunicodechar{⦄}{\ensuremath{|\!\}}}
\newunicodechar{𝓝}{\ensuremath{\mathcal{N}}}
\newunicodechar{≫}{\ensuremath{\gg}}
\newunicodechar{≪}{\ensuremath{\ll}}

\theoremstyle{plain}
\newtheorem{theorem}{Theorem}
\newtheorem{lemma}[theorem]{Lemma}
\newtheorem{proposition}[theorem]{Proposition}
\newtheorem{corollary}[theorem]{Corollary}
\newtheorem{conjecture}[theorem]{Conjecture}
\theoremstyle{definition}
\newtheorem{definition}[theorem]{Definition}
\newtheorem{remark}[theorem]{Remark}

\emergencystretch=3em
\raggedbottom
% Lean identifiers in running text (breakable at dots and underscores).
\expandafter\def\expandafter\UrlBreaks\expandafter{\UrlBreaks\do\_\do\.}
\DeclareUrlCommand\decl{\urlstyle{tt}}
\newcommand{\N}{\mathbb{N}}
\newcommand{\Stir}[2]{S(#1,#2)}
\newcommand{\nC}{n_{\mathrm C}}
\newcommand{\nN}{n_{\mathrm N}}
\newcommand{\nM}{n_{\mathrm M}}
\newcommand{\cnt}{\operatorname{cnt}}
\newcommand{\ins}{\operatorname{ins}}

\title[Partitions with $k$ blocks of mixed parity]{Set partitions of $\{1,\dots,n\}$ with $k$ blocks of mixed parity: a formal proof of Hanna's product formula (OEIS A124418, A124526)}
\author{math0}
\address{Produced by Claude (an AI model developed by Anthropic) in the math0 project}
\date{October 21, 2026}

\begin{document}
\begin{abstract}
Let $M(n,k)$ be the number of set partitions of $\{1,\dots,n\}$ with exactly $k$ blocks that contain both an odd and an even entry (OEIS A124418).
In 2006 P.~D.~Hanna conjectured that $M(n,k)=k!\,a(\lfloor n/2\rfloor,k)\,a(\lceil n/2\rceil,k)$, where $a(m,k)=\sum_i S(m,i)\binom ik$ (OEIS A049020) counts the partitions of an $m$-set with $k$ distinguished blocks; both OEIS entries still record the formula as a conjecture.
We prove it for all $n\ge0$ and $0\le k\le\lfloor n/2\rfloor$, the full index range of both triangles, by counting, by induction on the ground set, the partitions with given numbers of blocks meeting the odd entries, the even entries, and both.
The formula is not new: with the odd entries as the top row and the even entries as the bottom row of a partition diagram, mixed blocks are the transversal blocks, and the formula is the known count of partition diagrams of rank $k$ (East and Gray for even $n$; Dolinka, \DJ ur\dj ev and East for all $n$); our contribution is the first formal proof.
The proof is verified in Lean~4 with Mathlib against the frozen, fidelity-reviewed statement that math0 formalized from the OEIS entry; its definitions are checked by the kernel against the entry's data.
\end{abstract}
\maketitle

\section{Introduction}

A \emph{partition} of a finite set $s$ is a set of nonempty, pairwise disjoint subsets of $s$ (the \emph{blocks}) whose union is $s$.
Call a block $b\subseteq\N$ \emph{mixed} if it contains an odd and an even number.
The OEIS triangle A124418 \cite{A124418}, contributed by Emeric Deutsch on October~31, 2006, has entries
\[
  M(n,k)=\#\{\text{partitions of }[n]=\{1,\dots,n\}\text{ with exactly $k$ mixed blocks}\},\qquad 0\le k\le\lfloor n/2\rfloor .
\]
For example, $M(4,0)=4$, $M(4,1)=9$ and $M(4,2)=2$, the last two partitions being $12|34$ and $14|23$.
Let $S(m,i)$ denote the Stirling numbers of the second kind and
\[
  a(m,k)=\sum_{i=0}^{m}S(m,i)\binom ik ,
\]
the triangle A049020 \cite{A049020}: $a(m,k)$ is the number of partitions of an $m$-set in which $k$ of the blocks have been distinguished.
On November~8, 2006, Paul D.~Hanna contributed the triangle A124526 \cite{A124526},
\[
  T(n,k)=a(\lfloor n/2\rfloor,k)\,a(\lfloor (n+1)/2\rfloor,k),
\]
together with the conjecture $M(n,k)=k!\,T(n,k)$, recorded in both entries (Section~\ref{sec:statement}).
We prove it.

\begin{theorem}\label{thm:main}
For all $n\ge0$ and all $k$ with $0\le k\le\lfloor n/2\rfloor$,
\[
  M(n,k)=k!\,T(n,k)=k!\;a\bigl(\lfloor n/2\rfloor,k\bigr)\;a\bigl(\lceil n/2\rceil,k\bigr).
\]
\end{theorem}

This is exactly the statement verified in Lean (Section~\ref{sec:statement}): Hanna's conjecture over the full index range of both triangles.
The proof does not use the bound $k\le\lfloor n/2\rfloor$. The identity also holds, trivially, for $k>\lfloor n/2\rfloor$, where both sides vanish: each mixed block contains one of the $\lfloor n/2\rfloor$ even entries, and $\binom ik=0$ for $i\le\lfloor n/2\rfloor<k$. This extension is not part of the formal statement.
Summing over $0\le k\le\lfloor n/2\rfloor$ gives $\sum_k k!\,T(n,k)=B_n$, the Bell number, which is another comment in A124526.

\subsection*{Status and prior work}
The OEIS entries A124526 (revision \#14, November~2, 2020) and A124418 (revision \#26, March~13, 2026), as read by us and by the math0 referees from a local mirror of the OEIS data, still call the formula a conjecture.
Partial or equivalent forms appear in the OEIS without a written proof: the column $k=0$, $M(n,0)=B_{\lfloor n/2\rfloor}B_{\lceil n/2\rceil}$, is a formula by A.~P.~Heinz (2013) in A124419 \cite{A124419}, and is immediate because a partition without mixed blocks is a pair of partitions of the odd and of the even entries; Heinz's 2013 alternative Maple program in A124418 computes the triangle as $k!$ times a product of two factors $\sum_i\binom gi S(i,k)B_{g-i}$, which equal $a(g,k)$.

The theorem is a special case of known enumerations of partition diagrams, a fact not recorded in the OEIS entries.
A morphism $p\to q$ of the partition category $\mathcal P$ is a partition of the disjoint union of a top row $\{1,\dots,p\}$ and a bottom row $\{1',\dots,q'\}$; its \emph{transversal} blocks are those meeting both rows, and its \emph{rank} is their number.
Putting the $\lceil n/2\rceil$ odd entries of $[n]$ in the top row and the $\lfloor n/2\rfloor$ even entries in the bottom row identifies partitions of $[n]$ with such morphisms, and mixed blocks with transversal blocks.
East and Gray \cite[proof of Theorem~7.5 and Remark~7.7 of the arXiv version]{EG17} show that the $\mathscr J$-class $J_r(\mathcal P_m)$ of the partition monoid, the set of rank-$r$ partitions of $[m]\cup[m]'$, has $\rho_{mr}^2\,r!$ elements with $\rho_{mr}=\sum_j S(m,j)\binom jr=a(m,r)$; with $m=n/2$ this is Theorem~\ref{thm:main} for even $n$.
Halverson and Jacobson \cite[Remark~3.9]{HJ20} (Remark~3.23 in arXiv:1808.08118v2) give the same count in their notation: the partition algebra $P_k$ has $|W^m_k|^2m!$ basis diagrams of rank $m$, where $|W^m_k|=a(k,m)$.
Dolinka, \DJ ur\dj ev and East \cite[Proposition~4.11 and Remark~4.12 of the arXiv version]{DDE21} treat the rectangular case: their results give $r!\,a(p,r)\,a(q,r)$ rank-$r$ morphisms $p\to q$ in $\mathcal P$, which with $(p,q)=(\lceil n/2\rceil,\lfloor n/2\rfloor)$ is Theorem~\ref{thm:main} for all $n$.
These arguments count the possible top rows and bottom rows with $r$ distinguished blocks and the $r!$ ways to match them (in semigroup language: $\mathscr R$-classes, $\mathscr L$-classes and the size of an $\mathscr H$-class), which is the bijective argument of Remark~\ref{rem:bij} below.
None of these works mentions A124418 or A124526.
The math0 referees also checked Yaqubi and Mirzavaziri (arXiv:2608.30097), Ehrenborg, Hedmark and Hettle (J.~Integer Seq.\ 2017, on blocks of odd and even size) and Shattuck (arXiv:1401.6588), and found them not relevant.
We therefore classify Theorem~\ref{thm:main} as a known result, and the Lean proof described here as, to our knowledge, its first formal proof.
That proof does not construct the bijection of Remark~\ref{rem:bij}: it counts partitions of an arbitrary finite set by three statistics relative to an arbitrary predicate (Proposition~\ref{prop:refined}), inserting one element at a time, and then specializes to the predicate ``odd''.

\section{Statement}\label{sec:statement}

\subsection*{The OEIS conjecture}
The COMMENTS section of A124526, an entry by Paul D.~Hanna dated November~8, 2006, contains verbatim
\begin{quote}
\small\texttt{A124418(n,k) = k!*T(n,k) (conjecture).}
\end{quote}
and the FORMULA section of A124418 contains verbatim (here \texttt{T(n,k)} is the entry of A124418)
\begin{quote}
\raggedright\small\texttt{Conjecture: T(n,k) = k!*A049020([n/2],k)*A049020([(n+1)/2],k) where A049020(n,k)=Sum\_\{i=0..n\} S2(n,i)*C(i,k) and S2(n,k)=(1/k!)*Sum\_\{j=0..k\} (-1)\^{}(k-j)*C(k,j)*j\^{}n (the Stirling numbers of 2nd kind). - \_Paul D.\ Hanna\_, Nov 08 2006}
\end{quote}
Both triangles have offset $0$ and row $n$ indexed by $0\le k\le\lfloor n/2\rfloor$; $[x]$ is the floor.

\subsection*{Lean statement}
The statement is \decl{Math0.Statements.A124526.statement} in the math0 statement file \path{lean/Math0/Statements/OEIS/A124526.lean}, which was reviewed for fidelity and frozen (SHA-256 \texttt{aa643c9b\ldots c47325}) before any proof attempt.
Its definitions and the statement read, verbatim (docstrings, the data-check theorems and the restatement \decl{A124418_eq_card} omitted):
\begin{lstlisting}
namespace Math0.Statements.A124526
open Finset

def A049020 (m k : ℕ) : ℕ :=
  ∑ i ∈ range (m + 1), Nat.stirlingSecond m i * Nat.choose i k

def T (n k : ℕ) : ℕ := A049020 (n / 2) k * A049020 ((n + 1) / 2) k

def a (i : ℕ) : ℕ :=
  ((List.range (i + 1)).flatMap fun n => (List.range (n / 2 + 1)).map (T n)).getD i 0

def IsMixed (b : Finset ℕ) : Prop := (∃ x ∈ b, Odd x) ∧ (∃ x ∈ b, Even x)

instance : DecidablePred IsMixed := fun b => by unfold IsMixed; infer_instance

def A124418 (n k : ℕ) : ℕ :=
  (univ.filter fun P : Finpartition (Icc 1 n) => (P.parts.filter IsMixed).card = k).card

def statement : Prop :=
  ∀ n k : ℕ, k ≤ n / 2 → A124418 n k = k.factorial * T n k
\end{lstlisting}

\subsection*{Encoding and check against the OEIS data}
\texttt{A049020 m k} is $a(m,k)$, with Mathlib's \decl{Nat.stirlingSecond}, defined by the recurrence $S(0,0)=1$, $S(0,i+1)=S(m+1,0)=0$, $S(m+1,i+1)=(i+1)S(m,i+1)+S(m,i)$; this agrees with the explicit alternating sum quoted in the entry.
In $\N$, \texttt{n / 2} and \texttt{(n + 1) / 2} are $\lfloor n/2\rfloor$ and $\lfloor(n+1)/2\rfloor=\lceil n/2\rceil$, so \texttt{T n k} is A124526's $T(n,k)$, and \texttt{a} is the triangle read by rows with offset~$0$.
A Mathlib \texttt{Finpartition (Icc 1 n)} is a finite set of nonempty, pairwise disjoint finite sets with union $\{1,\dots,n\}$, and two such are equal iff they have the same parts; so \texttt{A124418 n k} is Deutsch's definition of $M(n,k)$, not the conjectured formula.
For $n=0$ the empty partition has no mixed block, matching row $0$ of A124418. The statement quantifies over the full index range of both triangles.

The statement file checks its definitions against the OEIS data by \texttt{decide +kernel}: \decl{T_data} and \decl{a_data} reproduce all $16$ listed rows and all $72$ listed terms of A124526; \decl{A124418_data} evaluates the \texttt{Finpartition}-based definition of A124418 on rows $n=0,\dots,3$, giving $1;\,1;\,1,1;\,2,3$ as listed (Mathlib enumerates \texttt{Finpartition}s through the power set of the power set of the ground set, and a kernel evaluation of row $n=4$ did not finish); and \decl{rhs_data} checks that $k!\,T(n,k)$ reproduces all $49$ listed terms of A124418 (rows $0,\dots,12$), which tests the transcription of the right-hand side but proves nothing about the conjecture.

\section{Proof}\label{sec:proof}

The proof follows the Lean file \path{lean/Math0/Frontier/OEIS/A124526.lean}; the names in parentheses are the corresponding Lean declarations.
Throughout, $c$ is a decidable predicate on $\N$, $s$ is a finite subset of $\N$, and partitions are as above.
Say that a block $b$ \emph{meets} $c$ if $c(x)$ holds for some $x\in b$.
For a partition $P$ of $s$ let $\nC(P)$, $\nN(P)$ and $\nM(P)$ be the numbers of blocks of $P$ that meet $c$, that meet $\neg c$, and that meet both $c$ and $\neg c$.
Let $\cnt_c(s;i,j,k)$ be the number of partitions $P$ of $s$ with $(\nC,\nN,\nM)(P)=(i,j,k)$, and
\[
  F(o,e,i,j,k)=S(o,i)\,S(e,j)\binom ik\binom jk k! .
\]

\begin{proposition}[\decl{cnt_eq}]\label{prop:refined}
Let $o=\#\{x\in s:c(x)\}$ and $e=\#\{x\in s:\neg c(x)\}$. Then $\cnt_c(s;i,j,k)=F(o,e,i,j,k)$ for all $i,j,k\ge0$.
\end{proposition}

\begin{remark}[Bijective meaning]\label{rem:bij}
$F(o,e,i,j,k)$ counts triples consisting of a partition of $s\cap c$ into $i$ blocks with $k$ of them marked, a partition of $s\setminus c$ into $j$ blocks with $k$ of them marked, and a bijection between the two sets of marked blocks.
Restricting a partition of $s$ to $s\cap c$ and to $s\setminus c$, marking the traces of the mixed blocks and matching the two traces of each mixed block is a bijection onto such triples, which proves Proposition~\ref{prop:refined} directly.
The formal proof uses the induction below instead, which needs only one simple bijection.
\end{remark}

\subsection{Inserting one element}
Let $t$ be finite, $a\notin t$, and $P$ a partition of $t$. For $x\in\{\varnothing\}\cup P$ define (\decl{insP}, \decl{ins})
\[
  \ins(P,\varnothing)=P\cup\bigl\{\{a\}\bigr\},\qquad
  \ins(P,b)=\bigl(P\setminus\{b\}\bigr)\cup\bigl\{b\cup\{a\}\bigr\}\quad(b\in P).
\]
Each is a partition of $t\cup\{a\}$: its blocks are nonempty subsets of $t\cup\{a\}$, and every $y\in t\cup\{a\}$ lies in exactly one of them (in Lean, \decl{Finpartition.ofExistsUnique} builds the partition from these three facts).

\begin{lemma}[\decl{ins_injective}, \decl{ins_surjective}, \decl{card_filter_ins}]\label{lem:ins}
The map $(P,x)\mapsto\ins(P,x)$, from pairs of a partition $P$ of $t$ and $x\in\{\varnothing\}\cup P$, to partitions of $t\cup\{a\}$, is a bijection. Consequently, for every property $f$ of partitions of $t\cup\{a\}$,
\[
  \#\{R: f(R)\}=\sum_{P}\#\bigl\{x\in\{\varnothing\}\cup P: f(\ins(P,x))\bigr\}.
\]
\end{lemma}

\begin{proof}
Since $a$ lies in no block of $P$, deleting $a$ from every block of $\ins(P,x)$ and discarding the empty set returns $P$ (\decl{rec_insP}); so $\ins(P,x)=\ins(P',x')$ implies $P=P'$.
Next, $\{a\}\in\ins(P,\varnothing)$ but $\{a\}\notin\ins(P,b)$, because $b\cup\{a\}\ne\{a\}$ ($b$ is nonempty) and no block of $P$ contains $a$ (\decl{insP_none_ne_some}).
Finally, if $\ins(P,b)=\ins(P,b')$ then $b\cup\{a\}$ is a block of $\ins(P,b')$ containing $a$, hence equals $b'\cup\{a\}$, and removing $a$ gives $b=b'$. This proves injectivity.
For surjectivity, let $R$ be a partition of $t\cup\{a\}$, let $B$ be its block containing $a$, and let $P$ be obtained from $R$ by deleting $a$ from every block and discarding the empty set (\decl{recP}); $P$ is a partition of $t$.
If $B=\{a\}$ then $R=\ins(P,\varnothing)$; otherwise $B\setminus\{a\}\in P$ and $R=\ins(P,B\setminus\{a\})$.
The counting formula is the cardinality of a set of pairs computed through this bijection (\decl{card_bij}).
\end{proof}

\begin{lemma}[\decl{card_filter_not_C}]\label{lem:notC}
The number of blocks of $P$ that do not meet $c$ is $\nN(P)-\nM(P)$.
\end{lemma}

\begin{proof}
The blocks meeting $\neg c$ split into the mixed ones and those not meeting $c$; conversely a block not meeting $c$ is nonempty and therefore meets $\neg c$.
\end{proof}

\begin{lemma}[\decl{stats_none}, \decl{stats_some_C}, \decl{stats_some_N}]\label{lem:stats}
Let $c(a)$ hold and let $P$ be a partition of $t$. Then
\begin{enumerate}
\item $(\nC,\nN,\nM)(\ins(P,\varnothing))=(\nC(P)+1,\nN(P),\nM(P))$;
\item if $b\in P$ meets $c$, then $(\nC,\nN,\nM)(\ins(P,b))=(\nC,\nN,\nM)(P)$;
\item if $b\in P$ does not meet $c$, then $(\nC,\nN,\nM)(\ins(P,b))=(\nC(P)+1,\nN(P),\nM(P)+1)$.
\end{enumerate}
\end{lemma}

\begin{proof}
(1) The new block $\{a\}$ meets $c$ but not $\neg c$.
In (2) and (3), the block $b$ is replaced by $b\cup\{a\}$, which meets $c$ and meets $\neg c$ iff $b$ does (as $c(a)$).
In (2), $b$ and $b\cup\{a\}$ therefore have the same three properties.
In (3), $b$ does not meet $c$, so it meets $\neg c$ (it is nonempty) and is not mixed, while $b\cup\{a\}$ meets $c$ and is mixed; so $\nC$ and $\nM$ increase by one and $\nN$ is unchanged.
\end{proof}

\begin{lemma}[\decl{card_choices}, \decl{cnt_insert}, \decl{sum1_zero}--\decl{sum3_succ}]\label{lem:rec}
Let $c(a)$ hold, and write $N(i,j,k)=\cnt_c(t;i,j,k)$ and $N'(i,j,k)=\cnt_c(t\cup\{a\};i,j,k)$. Then for all $i,j,k$,
\begin{align*}
  N'(0,j,k)&=0,\qquad N'(i+1,j,0)=N(i,j,0)+(i+1)N(i+1,j,0),\\
  N'(i+1,j,k+1)&=N(i,j,k+1)+(i+1)N(i+1,j,k+1)+(j-k)N(i,j,k),
\end{align*}
where $j-k$ is truncated subtraction in $\N$ (it only multiplies terms with $k\le j$).
\end{lemma}

\begin{proof}
Apply Lemma~\ref{lem:ins} with $f(R)$: $(\nC,\nN,\nM)(R)=(I,J,K)$. For fixed $P$ with statistics $(i',j',k')$, Lemmas~\ref{lem:stats} and~\ref{lem:notC} show that the choice $x=\varnothing$ gives $(i'+1,j',k')$, each of the $i'$ blocks meeting $c$ gives $(i',j',k')$, and each of the $j'-k'$ blocks not meeting $c$ gives $(i'+1,j',k'+1)$.
Summing over $P$, the first kind contributes $0$ if $I=0$ and $N(I-1,J,K)$ otherwise; the second contributes $I\cdot N(I,J,K)$, which is $0$ if $I=0$; the third contributes $0$ if $I=0$ or $K=0$, and $\bigl(J-(K-1)\bigr)\,N(I-1,J,K-1)$ otherwise.
\end{proof}

\begin{lemma}[\decl{step}]\label{lem:step}
$F(o+1,e,\cdot,\cdot,\cdot)$ is obtained from $F(o,e,\cdot,\cdot,\cdot)$ by the recurrence of Lemma~\ref{lem:rec}. Hence, if $c(a)$ holds and $\cnt_c(t;i,j,k)=F(o,e,i,j,k)$ for all $i,j,k$, where $o,e$ are the counts for $t$, then $\cnt_c(t\cup\{a\};i,j,k)=F(o+1,e,i,j,k)$ for all $i,j,k$.
\end{lemma}

\begin{proof}
Write $S=S(o,\cdot)$, $S'=S(e,\cdot)$ and use $S(o+1,0)=0$ and $S(o+1,i+1)=(i+1)S(o,i+1)+S(o,i)$.
If $i=0$, then $F(o+1,e,0,j,k)=0$.
If $k=0$, then $F(o+1,e,i+1,j,0)=\bigl((i+1)S(i+1)+S(i)\bigr)S'(j)$, which is the right-hand side.
For the indices $(i+1,j,k+1)$, expanding $S(o+1,i+1)$ and using Pascal's rule $\binom{i+1}{k+1}=\binom ik+\binom i{k+1}$, the difference between $F(o+1,e,i+1,j,k+1)$ and the right-hand side
$F(o,e,i,j,k+1)+(i+1)F(o,e,i+1,j,k+1)+(j-k)F(o,e,i,j,k)$ equals
\[
  S(i)\,S'(j)\binom ik k!\Bigl(\binom j{k+1}(k+1)-\binom jk(j-k)\Bigr)=0,
\]
by $(k+1)!=(k+1)k!$ and the identity $\binom j{k+1}(k+1)=\binom jk(j-k)$ (\decl{Nat.choose_succ_right_eq}, valid with truncated subtraction).
In Lean the last case is closed by \texttt{linear\_combination} over $\mathbb Z$ after this identity is cast.
\end{proof}

\begin{proof}[Proof of Proposition~\ref{prop:refined}]
Induction on $s$ (\decl{Finset.induction_on}), for all $(i,j,k)$ simultaneously.
For $s=\varnothing$ the only partition is the empty one, so $\cnt(\varnothing;i,j,k)=[i=j=k=0]$, and $F(0,0,i,j,k)=[i=0][j=0]\binom 0k^2k!=[i=j=k=0]$.
For $s=t\cup\{a\}$ with $a\notin t$: if $c(a)$, then $s$ has $o+1$ elements satisfying $c$ and $e$ satisfying $\neg c$, and Lemma~\ref{lem:step} applies.
If $\neg c(a)$, replace $c$ by $\neg c$. This exchanges $\nC$ and $\nN$ and fixes $\nM$, so $\cnt_{\neg c}(s;i,j,k)=\cnt_c(s;j,i,k)$ (\decl{cnt_not}), and $F(o,e,i,j,k)=F(e,o,j,i,k)$ (\decl{F_comm}).
The induction hypothesis for $t$ thus gives $\cnt_{\neg c}(t;\cdot)=F(e,o,\cdot)$; since $\neg c(a)$, Lemma~\ref{lem:step} for $\neg c$ gives $\cnt_{\neg c}(s;j,i,k)=F(e+1,o,j,i,k)$, that is, $\cnt_c(s;i,j,k)=F(o,e+1,i,j,k)$.
\end{proof}

\begin{proof}[Proof of Theorem~\ref{thm:main} (\decl{result})]
Take $c=\mathrm{Odd}$ and $s=[n]$. By induction on $n$, $[n]$ has $o=\lceil n/2\rceil=\lfloor(n+1)/2\rfloor$ odd and $e=\lfloor n/2\rfloor$ non-odd elements (\decl{card_odd}).
A block is \texttt{IsMixed} iff it meets $\mathrm{Odd}$ and $\neg\mathrm{Odd}$, since a natural number is even iff it is not odd; so $M(n,k)$ counts the partitions with $\nM=k$.
Each statistic is at most the number of blocks, which is at most $n$ (\decl{Finpartition.card_parts_le_card}).
Grouping the partitions with $\nM=k$ by $(\nC,\nN)\in\{0,\dots,n\}^2$ (\decl{card_eq_sum_card_fiberwise}) and applying Proposition~\ref{prop:refined},
\[
  M(n,k)=\sum_{i=0}^{n}\sum_{j=0}^{n}S(o,i)S(e,j)\binom ik\binom jk k!
  =k!\Bigl(\sum_{i=0}^{n}S(o,i)\binom ik\Bigr)\Bigl(\sum_{j=0}^{n}S(e,j)\binom jk\Bigr).
\]
Since $S(m,i)=0$ for $i>m$ and $o,e\le n$, the two sums are $a(o,k)$ and $a(e,k)$ (\decl{sum_eq_A049020}), so $M(n,k)=k!\,a(e,k)\,a(o,k)=k!\,T(n,k)$.
The hypothesis $k\le\lfloor n/2\rfloor$ of \decl{statement} is introduced and not used.
\end{proof}

\section{Machine verification}\label{sec:lean}

\subsection*{Files and declarations}
The proof is the file \path{lean/Math0/Frontier/OEIS/A124526.lean} of the math0 repository, \url{https://github.com/Foundation-0/math0}.
It imports the frozen statement file \path{lean/Math0/Statements/OEIS/A124526.lean} and the audit commands in \path{lean/Math0/Audit.lean}.
The main declaration is \decl{Math0.Frontier.A124526.result}, declared as \texttt{theorem result : statement}; its type is the constant \decl{Math0.Statements.A124526.statement}.
All other declarations are private; the main ones are named in Section~\ref{sec:proof}.
The private definitions \texttt{nC}, \texttt{nN}, \texttt{nM}, \texttt{cnt} and \texttt{F} are those of Section~\ref{sec:proof}, with \texttt{Finpartition s} for partitions and \texttt{\#(P.parts.filter $\cdots$)} for the numbers of blocks; Proposition~\ref{prop:refined} reads
\begin{lstlisting}
private theorem cnt_eq (s : Finset ℕ) :
    ∀ i j k, cnt c s i j k = F #(s.filter c) #(s.filter fun x => ¬ c x) i j k
\end{lstlisting}
No bounded computation is used, and the file has no \texttt{sorry}, \texttt{native\_decide}, \texttt{set\_option}, axioms or custom syntax.



\subsection*{Audit}
The file ends with the command \texttt{\#math0\_settles}~\decl{Math0.Frontier.A124526.result}~\decl{Math0.Statements.A124526.statement}.
This command, defined in \path{lean/Math0/Audit.lean} for statements that math0 formalizes itself, fails unless (1) \decl{statement} is a universe-monomorphic \texttt{Prop}-valued definition; (2) \decl{statement} and \decl{result} use no axioms besides \decl{propext}, \decl{Classical.choice} and \decl{Quot.sound} (this rejects \decl{sorryAx}, the axiom \decl{Lean.ofReduceBool} behind \texttt{native\_decide}, and custom axioms); and (3) the type of \decl{result} is syntactically \decl{statement} or its negation, without unfolding.
Running \texttt{lake env lean Math0/Frontier/OEIS/A124526.lean} in the directory \path{lean} prints\par\nobreak\noindent
\begin{minipage}{\linewidth}
\begin{lstlisting}[basicstyle=\ttfamily\scriptsize]
math0 audit: 'Math0.Statements.A124526.statement' depends on axioms [Classical.choice, Quot.sound, propext]
math0 audit: 'Math0.Frontier.A124526.result' depends on axioms [Classical.choice, Quot.sound, propext]
math0 statement check: 'Math0.Frontier.A124526.result' proves 'Math0.Statements.A124526.statement'
\end{lstlisting}
\end{minipage}
so \decl{result} proves the statement using only the three permitted axioms.

\subsection*{Toolchain and reproduction}
We use Lean~4 \cite{Lean4} \texttt{v4.33.1} and Mathlib \cite{Mathlib} \texttt{v4.33.1}, pinned in \path{lean/}; the repository also pins formal-conjectures \cite{FC} at commit \texttt{c77fdd4}, which this result does not use (the statement is math0's own).
To reproduce, run \path{scripts/verify.sh} from the repository root: it fetches the Mathlib cache, runs \texttt{lake build} (re-running every audit command), and checks that the pinned dependencies are unmodified.
For this result alone, run \texttt{lake build Math0.Statements.OEIS.A124526} and then the command above in \path{lean}; the proof file elaborates in under half a minute.

{\small
\begin{thebibliography}{99}

\bibitem{DDE21}
I.~Dolinka, I.~\DJ ur\dj ev and J.~East, \emph{Sandwich semigroups in diagram categories}, Internat.\ J.\ Algebra Comput.\ \textbf{31} (2021), no.~7, 1339--1404, doi:10.1142/S021819672150048X; arXiv:1910.10286.

\bibitem{EG17}
J.~East and R.~D.~Gray, \emph{Diagram monoids and Graham--Houghton graphs: idempotents and generating sets of ideals}, J.\ Combin.\ Theory Ser.~A \textbf{146} (2017), 63--128, doi:10.1016/j.jcta.2016.09.001; arXiv:1404.2359.

\bibitem{FC}
M.~Firsching, P.~Lezeau, S.~Mercuri, M.~Z. Horv\'ath, Y.~Dillies, C.~S\"onne, E.~Wieser, F.~Zhang, T.~Hubert, B.~Ag\"uera y Arcas, and P.~Kohli, \emph{Formal Conjectures: An Open and Evolving Benchmark for Verified Discovery in Mathematics}, arXiv:2605.13171, 2026. Repository: \url{https://github.com/google-deepmind/formal-conjectures}.

\bibitem{HJ20}
T.~Halverson and T.~N.~Jacobson, \emph{Set-partition tableaux and representations of diagram algebras}, Algebr.\ Comb.\ \textbf{3} (2020), no.~2, 509--538, doi:10.5802/alco.102; arXiv:1808.08118.

\bibitem{Lean4}
L.~de Moura and S.~Ullrich, \emph{The Lean 4 theorem prover and programming language}, in: Automated Deduction -- CADE 28, Lecture Notes in Comput. Sci. \textbf{12699}, Springer, 2021, 625--635.

\bibitem{Mathlib}
The mathlib Community, \emph{The Lean mathematical library}, in: Proceedings of the 9th ACM SIGPLAN International Conference on Certified Programs and Proofs (CPP 2020), ACM, 2020, 367--381.

\bibitem{A049020}
The OEIS Foundation Inc., \url{https://oeis.org/A049020}.

\bibitem{A124418}
The OEIS Foundation Inc., \url{https://oeis.org/A124418}, entry by E.~Deutsch, Oct.~31, 2006.

\bibitem{A124419}
The OEIS Foundation Inc., \url{https://oeis.org/A124419}, entry by E.~Deutsch, Oct.~31, 2006.

\bibitem{A124526}
The OEIS Foundation Inc., \url{https://oeis.org/A124526}, entry by P.~D.~Hanna, Nov.~8, 2006.

\end{thebibliography}
}

\end{document}
